\section{Limitations and Future Work}
\label{sec:limitations}

The theorem proved by the artifact is a coordinate-represented compact-support
Stokes theorem.  This is a strong local-to-global result, but it is not the
final possible theorem about native manifold integrals.  This section states
the boundary of the claim and the natural next steps.

\subsection{Represented-to-native comparison}

The main remaining mathematical bridge is to identify the generated represented
endpoints with a chosen native integration API.  The desired comparison would
identify \code{canonicalRepresentedBulkIntegral} with \(\int_M\dd\omega\) and
\code{canonicalRepresentedBoundaryIntegral} with
\(\int_{\partial M}\omega\).
Such a bridge should prove that the finite coordinate sums generated by the
local proof are independent of the selected cover and partition and agree with
the native integral definitions.  This requires a stable target API for
integration of differential forms on oriented smooth manifolds with boundary.

The current theorem deliberately stops before this comparison.  It proves the
finite coordinate equality generated by Stokes, leaving the comparison with
native integrals as a separate theorem layer.

\subsection{Orientation and boundary APIs}

The half-space layer already proves the local sign bridge between the lower
normal face and the outward-first convention.  A full native theorem will need
to connect this local sign convention to a global boundary orientation API.
That includes transporting orientations through charts, proving compatibility
under chart changes, and connecting boundary measures or integrals with the
coordinate lower-face terms.

\subsection{Consolidation}

The repository still contains historical routes and compatibility modules.
They are valuable as development records and as local theorem sources, but the
public API should remain centered on the first-principles route.  Future
engineering work can continue to consolidate implementation modules, move
historical theorem shapes out of default imports, and keep theorem-facing
names focused on the final mathematical interface.

\subsection{Possible extensions}

Several extensions are natural:
\begin{itemize}
  \item comparison of represented endpoints with native integrals;
  \item a compact oriented smooth manifold with boundary theorem once the
  native integration bridge is in place;
  \item more systematic chart-independence lemmas for represented endpoints;
  \item upstreaming general-purpose half-space and support-localization lemmas
  to mathlib where appropriate;
  \item a shorter conference version emphasizing the artifact architecture and
  verification gates.
\end{itemize}

These are extensions of the current milestone, not repairs of it.  The present
artifact already proves the generated coordinate Stokes equality from
first-principles compact-support hypotheses.
